\section{Preliminary Feasibility Evaluation}
\label{sec:preliminary}

We conducted a developer-assisted feasibility evaluation using a simulated
stock-trading multi-agent system. The system comprises market-data,
analysis, portfolio, risk, execution, and reconciliation agents communicating
through a FIFO message queue in a single process. Market inputs and exchange
faults are synthetic; the experiment does not involve live trading.

\paragraph{Artifacts and provenance.}
A conversational coding assistant authored six restricted Python policies
covering timeout retries, quote freshness, cash reservation, partial-fill
accounting, execution-notification deduplication, and late-fill reconciliation.
Three further configurations combined these policies to update two, two,
and four agent policies simultaneously. These configurations reuse the
single-policy artifacts and are not independent generations. The assistant
had previously developed and inspected the evaluator, reference policies,
and tests. Consequently, this experiment evaluates the feasibility of
integrating these particular artifacts rather than independent LLM
performance on unseen tasks. No separate API calls were used, and an exact
backend model snapshot was unavailable.

\paragraph{Procedure.}
The candidate sources were frozen and hashed before execution. The tool
checked their supported Python syntax and command schemas, then integrated
them into the simulator. Development and final assessment suites covered
ordinary operation, request loss, lost acknowledgements, pending-order cash
contention, partial fills, duplicate notifications, cancellation, and
interleaved operations. A combined recovery scenario exercised several
updated policies within one transaction trace. Each scenario family was
executed with twenty input seeds. These seeds represent test executions,
not independent LLM generations. Although the suites have development and
final roles in the tool, the final suite was not unseen by the assistant
in this developer-assisted run.

\begin{table}[t]
\centering
\caption{Observed acceptance of the frozen artifacts. Counts denote scenario
executions, not independent generations.}
\label{tab:preliminary-results}
\begin{tabular}{lrr}
\hline
Configuration & Development & Final \\
\hline
Six individual policies, combined & 600/600 & 1800/1800 \\
Market-data and execution bundle & 180/180 & 380/380 \\
Risk and execution bundle & 180/180 & 380/380 \\
Four-policy bundle & 260/260 & 460/460 \\
\hline
Total & 1220/1220 & 3020/3020 \\
\hline
\end{tabular}
\end{table}

\paragraph{Results.}
All six individual policy artifacts and the three derived configurations
passed their assessment suites, yielding 4,240 candidate scenario executions
without observed violations. The corresponding pre-evolution baselines,
which intentionally lacked the requested capabilities, were rejected by
the final assessments. Including baseline executions, the experiment
comprised 8,480 scenario executions. No repair attempts were observed because
all submitted candidates passed on their first assessment. Thus, the results
provide no empirical estimate of feedback-guided repair effectiveness or
an advantage over independent resampling.

\paragraph{Limitations.}
The evaluator-aware generation process, small handcrafted policies,
shared simulator, overlapping requirements, and reuse of policies across
bundles limit the conclusions. Passing the selected executions is not a
formal correctness proof. The study does not establish generalization,
distributed-system correctness, profitability, or a population-level LLM
success rate. Independent generation sessions, additional systems, and
budget-controlled repair comparisons are required for stronger claims.
